% Generated draft: requires mathematical and source review.
\documentclass[11pt]{article}
\usepackage[margin=25mm]{geometry}
\usepackage{fontspec}
\setmainfont{Latin Modern Roman}
\usepackage{amsmath,amssymb,mathtools,unicode-math}
\setmathfont{Latin Modern Math}
\usepackage{enumitem,needspace,xurl,hyperref}
\setlength{\parindent}{0pt}
\setlength{\parskip}{5pt}
\setlength{\emergencystretch}{3em}
\newcommand{\N}{\mathbb{N}}
\newcommand{\Z}{\mathbb{Z}}
\newcommand{\Q}{\mathbb{Q}}
\newcommand{\R}{\mathbb{R}}
\newcommand{\C}{\mathbb{C}}
\newcommand{\F}{\mathbb{F}}
\newcommand{\A}{\mathbb{A}}
\newcommand{\PP}{\mathbb{P}}
\newcommand{\E}{\mathbb{E}}
\newcommand{\eps}{\varepsilon}
\newcommand{\dd}{\,\mathrm{d}}
\begin{document}
\section*{UnsolvedMath ID: 30004500}
\textbf{Source problem code:} catalogue entry
\subsection{Definitions and mathematical statement}
Put \(k=\mathbb F_q(t)\), where \(q\) is a fixed prime power, and fix a quaternion algebra \(D\) over \(k\). For a quadratic extension \(F/k\), write \(D_F=D\otimes_k F\), and let \(g\) denote the genus of \(F\). For a global function field \(L\), let \(\mathbb A_L\) denote its adele ring. Write \(G_k=D^\times\), \(G_F=D_F^\times\), and let \(Z_F\) be the center of \(G_F\), so that \(Z_F(\mathbb A_F)=\mathbb A_F^\times\). The extension \(F/k\) induces the inclusion of adelic groups used below.

Let \(K\subseteq G_F(\mathbb A_F)\) be a compact subgroup, and put
\[
H_K=K\cap G_k(\mathbb A_k).
\]
The supplied record defines
\[
m=[M:H_K],
\]
where \(M\) is a maximal compact subgroup of \(G_k(\mathbb A_k)\) containing \(H_K\). The choice of \(M\) and the finiteness of this index are not specified. In particular, compactness of \(K\) alone does not supply the finite-index hypothesis needed to use \(m\) as a positive integer parameter.

Consider a normalized Hecke eigenfunction \(f\) on
\[
X_{F,K}=G_F(F)\backslash G_F(\mathbb A_F)\big/\bigl(KZ_F(\mathbb A_F)\bigr).
\]
The source does not specify the normalization, the Hecke operators, or the automorphic function space from which \(f\) is taken. Its proposed period is
\[
P_F(f)=\int_{G_k(k)\backslash G_k(\mathbb A_k)}f(x)\,\mathrm d\mu(x),
\]
where \(f\) is restricted along the adelic inclusion. The measure \(\mu\), its normalization, and convergence of this integral are not specified. The displayed integration domain has not been divided by the center; it is retained exactly as supplied.

\textbf{Supplied lower-bound question.} With \(q\) and \(D\) fixed, does there exist a constant \(C\), depending only on \(q\) and \(D\), such that, as \(F/k\) varies over quadratic extensions of genus \(g\), \(K\) varies over the stated compact subgroups, and \(f\) varies over every normalized Hecke eigenfunction on \(X_{F,K}\), one has
\[
P_F(f)\stackrel{?}{=}\Omega\!\left(m^{-C}q^{g/2}\right)?
\]
This is the supplied target, with unresolved qualifications. In particular, a complex-valued period has no specified ordered lower bound, and the source does not state whether an absolute value is intended. Under an absolute-value interpretation with fully uniform implied constants, the proposed inequality would take the form
\[
|P_F(f)|\ge c\,m^{-C}q^{g/2},
\]
with \(c>0\) independent of \(F,K,f\). That interpretation is not adopted as a correction: the intended phase convention or absolute value, the dependence of \(c\), and whether the bound applies for every genus or only beyond a uniform genus threshold remain unspecified. The source also does not impose a nonvanishing or distinction condition on \(f\). These issues prevent an unqualified precise inequality from being extracted.

The supplied dated curation describes this as open and says that the source passage checks only the split standard-level case, while the proposed universal bound over varying extensions and compact subgroups remains unproved. The record does not specify that special case sufficiently to formulate a separate verified assertion.
\subsection{Short English statement}
For a fixed quaternion algebra over \(\mathbb F_q(t)\), the supplied question asks for a uniform lower bound of size \(m^{-C}q^{g/2}\) for periods of every normalized Hecke eigenfunction over varying quadratic extensions. A precise formulation still requires the intended absolute value or phase convention, normalization of functions and measures, convergence of the period, a finite level index, and the scope of the asymptotic quantifiers. The supplied open status and split standard-level qualification are dated curation requiring independent review.
\subsection{Sources}
\begin{itemize}
\item Local UnsolvedMath registry, \texttt{data/unsolvedmath/problems.json}: permanent ID \(30004500\), title ``Lower Bounds for Function-Field Quaternionic Periods,'' source problem code \texttt{catalogue entry}, in the Oberwolfach Reports Open Problems collection. This draft uses the supplied complete registry record.
\item Supplied report citation: \emph{Automorphic Forms and Arithmetic} (2021), Oberwolfach Reports, pp.\ 1317--1383, DOI \texttt{10.4171/owr/2020/26}. Supplied source URL and literature-source link, labeled ``Oberwolfach Report 26/2020'': \url{https://doi.org/10.4171/owr/2020/26}. The record gives proposed year 2020 and does not identify an individual proposer.
\item Literature assessment dated exactly \texttt{2026-08-22}: the supplied status is ``open.'' The supplied assessment says that the function-field lower-bound question is checked only in the split standard-level case in the source passage, whereas the proposed universal bound over varying extensions and compact subgroups remains unproved. This is supplied dated curation, not an independently verified literature assessment; its mathematical claims and time-sensitive status require independent verification.
\item The registry labels its statement \texttt{corrected\_verified} and reports that the official report was checked to restore the quaternion algebra, varying quadratic extension, genus, compact subgroup, index, automorphic quotient, and parameter dependence. This is a supplied verification claim, not verification performed for this draft. The supplied citation uses the year 2021, while the report label and DOI use 2020; both are retained.
\end{itemize}
\end{document}
